\documentclass[10pt]{article}
\usepackage[T1]{fontenc}
\usepackage[utf8]{inputenc}
\usepackage[french]{babel}
\usepackage{amsmath,amssymb,mathrsfs,mathpazo}
\usepackage{geometry}
\geometry{a4paper,margin=2cm}
\usepackage{graphicx,array,tabularx,tabularray,colortbl,multirow,multicol,booktabs,enumitem,pifont,ulem,fancybox,diagbox,cancel}
\usepackage{pstricks,pstricks-add,pst-plot,pst-tree,pst-node,pst-text,pst-func,pst-math,pst-bspline,pst-3dplot,pst-eucl,pst-solides3d}
\let\pscalculate\undefined
\usepackage{pst-bezier}
\usepackage{tikz,pgfplots,tkz-tab}
\usetikzlibrary{arrows,arrows.meta,patterns,calc,positioning,trees,decorations.markings,shapes.geometric,shapes.misc,shapes.arrows,intersections,decorations.pathmorphing,decorations.pathreplacing,backgrounds,shadows,babel}
\pgfplotsset{compat=1.18}
\usepgfplotslibrary{fillbetween}
\usepackage[np]{numprint}
\usepackage{listings,makecell,pas-tableur,twoopt}
\let\barre\undefined
\usepackage[pstricks]{tablvar}
\definecolor{darkgreen}{RGB}{0,100,0}
\providecommand{\red}{\color{red}}
\providecommand{\blue}{\color{blue}}
\providecommand{\green}{\color{green}}
\providecommand{\black}{\color{black}}
\providecommand{\white}{\color{white}}
\providecommand{\gray}{\color{gray}}
\providecommand{\cyan}{\color{cyan}}
\providecommand{\magenta}{\color{magenta}}
\providecommand{\yellow}{\color{yellow}}
\providecommand{\violet}{\color{violet}}
\providecommand{\orange}{\color{orange}}
\providecommand{\euro}{\texteuro}
\providecommand{\EUR}{\texteuro}
\providecommand{\ieme}{\textsuperscript{e}}
\providecommand{\degres}{\ensuremath{^\circ}}
\providecommand{\degree}{\ensuremath{^\circ}}
\providecommand{\vv}[1]{\overrightarrow{#1}}
\providecommand{\Coord}[1]{#1}
\providecommand{\decosix}{}
\providecommand{\starredbullet}{\ensuremath{\bullet}}
\pagestyle{empty}
\setlength{\parindent}{0pt}
\usepackage{adjustbox}\DeclareUnicodeCharacter{2236}{:}\DeclareUnicodeCharacter{2190}{\ensuremath{\leftarrow}}
\providecommand{\R}{\mathbb{R}}
\providecommand{\N}{\mathbb{N}}
\providecommand{\Z}{\mathbb{Z}}
\providecommand{\Q}{\mathbb{Q}}
\providecommand{\D}{\mathbb{D}}
\providecommand{\C}{\mathbb{C}}
\providecommand{\e}{{\mathrm{e}}}
\providecommand{\vect}[1]{\overrightarrow{#1}}
\providecommand{\vectt}[1]{\overrightarrow{\mathrm{#1}}}
\providecommand{\vvt}[1]{\overrightarrow{\mathrm{#1}}}
\providecommand{\np}[2][]{{#2~#1}}
\providecommand{\vv}[1]{{\overrightarrow{#1}}}
\renewcommand{\d}{{\mathrm{d}}}
\providecommand{\degres}{{^\circ}}
\providecommand{\degre}{{^\circ}}
\providecommand{\blue}{\color{blue}}
\providecommand{\red}{\color{red}}
\providecommand{\enskip}{\quad}
\providecommand{\llbracket}{\lbrack\!\lbrack}
\providecommand{\rrbracket}{\rbrack\!\rbrack}
\providecommand{\euro}{{\text{€}}}
\providecommand{\no}{{\text{n°}}}
\providecommand{\interval}[2]{{[#1;#2]}}
\providecommand{\intervFO}[2]{{[#1;#2[}}
\providecommand{\dotfill}{\ldots}

\providecommand{\ds}{\displaystyle}
\providecommand{\textding}[1]{#1}
\providecommand{\Oij}{\(\mathrm{O};\vec{\imath},\vec{\jmath}\)}
\providecommand{\Oijk}{\(\mathrm{O};\vec{\imath},\vec{\jmath},\vec{k}\)}
\providecommand{\Ouv}{\(\mathrm{O};\vec{u},\vec{v}\)}
\providecommand{\intervFF}[2]{\ensuremath{[#1;#2]}}
\renewcommand{\euro}{\ifmmode\text{\texteuro}\else\texteuro\fi}
\DeclareUnicodeCharacter{00B0}{\ifmmode^{\circ}\else\textdegree\fi}
\DeclareUnicodeCharacter{20AC}{\ifmmode\text{\texteuro}\else\texteuro\fi}
% Énoncés : Éducation nationale. Transcription et corrigés : APMEP, Denis Vergès et contributeurs. Découpage et classement : CoopMaths, Sébastien Lozano, Éric Elter et contributeurs. Adaptation de la présentation : Maths Explorer.
% CC BY-SA : https://coopmaths.fr/www/about/
% https://coopmaths.fr/alea/static/bac/2024/tex/bac_2024_06_sujet1_suede_4.tex
\begin{document}
\noindent\textit{APMEP / CoopMaths. CC BY-SA. Présentation adaptée par Maths Explorer.}\par\medskip
On modélise un passage de spectacle de voltige aérienne en duo de la manière suivante :

\begin{itemize}[label=$\bullet$]
\item on se place dans un repère orthonormé \Oijk, une unité représentant un mètre;

\item l'avion n\up{o} 1 doit relier le point O au point $\mathrm{A}(0~;~200~;~0)$ selon une trajectoire rectiligne, à la vitesse constante de $\np[m/s]{200}$;

\item l'avion n\up{o} 2 doit, quant à lui, relier le point $\mathrm{B}(-33~;~75~;~44)$ au point $\mathrm{C}(87~;~75~;~-116)$ également selon une trajectoire rectiligne, et à la vitesse constante de $\np[m/s]{200}$.

\item au même instant, l'avion \no 1 est au point O et l'avion \no 2 est au point B.
\end{itemize}
\begin{tikzpicture}[x={(-160:0.4mm)},y={(0:0.35mm)},z={(90:0.3mm)},>=stealth]
	\draw[->] (-150,0,0) -- (150,0,0) node[below left]{$x$};
	\draw[->] (0,0,0) -- (0,250,0) node[below]{$y$};
	\draw[->] (0,0,-200) -- (0,0,80) node[left]{$z$};
	\fill (0,0,0) circle (3pt) node[above right]{O};
	\fill (0,200,0) circle (3pt) node[above right]{A};
	\fill (87,75,-116) circle (3pt) node[below=5pt]{C};
	\fill (-33,75,44) circle (3pt) node[above right]{B};
	\draw[line width=2pt] (0,0,0) -- (0,200,0) (87,75,-116) -- (-33,75,44) ;
	\draw[teal,dashed,line width=1pt] (0,0,0) -- (0,75,0) node[color=black,below right]{75} -- (-33,75,0) -- (-33,75,44) -- (-33,0,44) -- (0,0,44)node[black,left]{44} -- cycle
	(0,0,44) -- (0,75,44) -- (-33,75,44) (0,75,44) -- (0,75,0)
	(0,0,0) -- (-33,0,0)node[black,above left]{$-33$} -- (-33,75,0) (-33,0,0) -- (-33,0,44);
	\draw[red,dashed,line width=1pt] (0,0,0) -- (0,75,0) -- (0,75,-116) -- (87,75,-116) -- (87,0,-116) -- (87,0,0)node[black,above left]{87} -- cycle
	(87,0,0)--(87,75,0) -- (87,75,-116) (87,75,0) -- (0,75,0)
	(0,0,0) -- (0,0,-116)	(87,0,-116) -- (0,0,-116)node[black, above left]{$-116$} --(0,75,-116);
\end{tikzpicture}
\begin{enumerate}
\item Justifier que l'avion \no 2 mettra autant de temps à parcourir le segment [BC] que l'avion \no 1 à parcourir le segment [OA].
\item Montrer que les trajectoires des deux avions se coupent.
\item Les deux avions risquent-ils de se percuter lors de ce passage ?
\end{enumerate}

\end{document}
